\section{Determinant coefficients and unconditional frequency dispersion}
\label{sec:tpc32-determinant-dispersion}

We now retain the three cutoff-shell channels in their matched signed
combination.  Write
\begin{equation}
 \cS_{\rm sh}(\cE)
 =\sum_{\alpha,\gamma,j}
   \mathfrak a^{\rm sh}_{\alpha,\gamma,j}
   \one_{\cE(\alpha,\gamma,j)},
 \label{eq:tpc32-matched-physical-sum}
\end{equation}
where $\mathfrak a^{\rm sh}_{\alpha,\gamma,j}$ is the complete signed
physical summand: it contains both opened-row coefficients, the actual
row-pair mask, the fixed-period and smooth factors, and the matched
kernel
\[
 \mathsf A_{m_\alpha,U_0}(j)\mathsf A_{m_\gamma,U_0}(j)
 -\mathsf A_{m_\alpha,T}(j)\mathsf A_{m_\gamma,T}(j).
\]
Thus no triangle inequality separates the two mixed channels from the
both-ultra channel in the definition below.

Put
\begin{equation}
 c_{\alpha,\gamma}(j)
 =(m_\alpha j+h,m_\gamma j+h),
 \qquad
 \Delta^\#_{\alpha,\gamma}(j)
 =\frac{m_\alpha-m_\gamma}{c_{\alpha,\gamma}(j)}.
 \label{eq:tpc32-dispersion-content-determinant}
\end{equation}
For $1\le C\ll_{\mathscr D}Q$, define the signed
normalized-determinant coefficient
\begin{equation}
 A_C(n)
 :=\cS_{\rm sh}
 \bigl(c_{\alpha,\gamma}(j)\le C,
       \ \Delta^\#_{\alpha,\gamma}(j)=n\bigr),
 \qquad n\in\Z.
 \label{eq:tpc32-determinant-coefficient}
\end{equation}
The off-diagonal row mask makes $A_C(0)=0$.  To avoid confusing this
coefficient value with a Fourier frequency, we shall always write
$\widehat A_{C,q}(0)$ for the latter.

The physical row interval has length $O_{\mathscr D}(Q)$.  Fix a
constant $B_{\mathscr D}\ge1$ such that
\begin{equation}
 \operatorname{supp}A_C
 \subset
 [-B_{\mathscr D}Q,B_{\mathscr D}Q]\cap
 (\Z\setminus\{0\})
 \label{eq:tpc32-determinant-support-interval}
\end{equation}
for every $C$.  Indeed,
$|\Delta^\#|=|m_\alpha-m_\gamma|/c\le
B_{\mathscr D}Q$.  In particular the diameter of this support, rather
than merely its cardinality, is $O_{\mathscr D}(Q)$; this is the
quantity entering the additive large sieve below.

\subsection{\texorpdfstring{From the inherited \(\ell^1\) and singleton
bounds to energy}{From the inherited l1 and singleton bounds to energy}}

Let $A_C^+(n)$ be the absolute majorant formed by summing the absolute
values of the complete matched physical summands in
\eqref{eq:tpc32-determinant-coefficient}.  The physical absolute-mass
interface and the singleton version of the determinant-support theorem
of TPC-31 give, for every fixed $\eps>0$,
\begin{align}
 \sum_n A_C^+(n)
 &\ll_{\eps,h,\mathscr D}X^\eps N_0,
 \qquad N_0:=JQ^2\asymp XQ,
 \label{eq:tpc32-inherited-det-l1}\\
 \max_n A_C^+(n)
 &\ll_{\eps,h,\mathscr D}
 X^\eps Q\log(2L)
 \{C+J\log(2C)\}.
 \label{eq:tpc32-inherited-det-linfty}
\end{align}
For clarity, the second estimate follows by taking the allowed
determinant set at each exact content to be the singleton $\{n\}$:
its physical support count is at most one, and
\[
 \sum_{c\le C}\left(1+\frac Jc\right)
 \ll C+J\log(2C).
\]
Both estimates remain valid for the matched sum because the absolute
majorant of the sum of its three expanded channels is bounded by the
sum of their TPC-31 majorants \citep{WangTPC31}.

\begin{proposition}[Determinant coefficient norms]
\label{prop:tpc32-determinant-norms}
For every $1\le C\ll_{\mathscr D}Q$,
\begin{align}
 \|A_C\|_1
 &\ll_{\eps,h,\mathscr D}X^\eps N_0,
 \label{eq:tpc32-det-l1}\\
 \|A_C\|_\infty
 &\ll_{\eps,h,\mathscr D}X^\eps Q(C+J),
 \label{eq:tpc32-det-linfty}\\
 \|A_C\|_2^2
 &\ll_{\eps,h,\mathscr D}
 X^\eps JQ^3(C+J).
 \label{eq:tpc32-det-l2-general}
\end{align}
If $C\le J$, then
\begin{equation}
 \boxed{
 \|A_C\|_2^2
 \ll_{\eps,h,\mathscr D}
 X^\eps\frac{N_0^2}{Q}.}
 \label{eq:tpc32-det-l2-C-below-J}
\end{equation}
\end{proposition}

\begin{proof}
Since $|A_C(n)|\le A_C^+(n)$,
\eqref{eq:tpc32-inherited-det-l1} gives
\eqref{eq:tpc32-det-l1}.  Equation
\eqref{eq:tpc32-inherited-det-linfty} gives
\eqref{eq:tpc32-det-linfty} after enlarging the soft exponent to absorb
$\log(2L)\log(2C)$.  Apply both inherited estimates with $\eps/2$ and
use
\[
 \sum_n|A_C(n)|^2
 \le\sum_n A_C^+(n)^2
 \le \bigl(\max_nA_C^+(n)\bigr)
       \sum_nA_C^+(n).
\]
This proves \eqref{eq:tpc32-det-l2-general}, after relabelling the
soft exponent.  If $C\le J$, then
\[
 JQ^3(C+J)\ll J^2Q^3
 =\frac{(JQ^2)^2}{Q}=\frac{N_0^2}{Q},
\]
which proves the last assertion.
\end{proof}

The argument is deliberately sign-blind only at the final norm bound.
The coefficient $A_C$ itself was formed after the three physical
channels were matched.  Thus every future improvement of
\eqref{eq:tpc32-det-l2-C-below-J} coming from their signed interaction
can be inserted without changing the Fourier statements below.

\subsection{A no-wrap determinant Fourier transform}

Choose an integer
\[
 Y>2B_{\mathscr D}Q+2.
\]
Bertrand's postulate supplies, for all sufficiently large $Q$, a prime
$q$ with $Y<q<2Y$.  We fix $Y$ to be the least integer above the
displayed threshold.  Consequently
\begin{equation}
 q\asymp_{\mathscr D}Q,
 \qquad
 q>\operatorname{diam}(\operatorname{supp}A_C).
 \label{eq:tpc32-no-wrap-prime}
\end{equation}
Define
\begin{equation}
 \widehat A_{C,q}(r)
 :=\sum_{n\in\Z}A_C(n)e_q(-rn),
 \qquad
 e_q(x)=\exp(2\pi ix/q),
 \qquad r\pmod q.
 \label{eq:tpc32-no-wrap-DFT}
\end{equation}

\begin{theorem}[No-wrap Parseval and exceptional frequencies]
\label{thm:tpc32-no-wrap-parseval}
Let $q$ satisfy \eqref{eq:tpc32-no-wrap-prime}.  Then
\begin{equation}
 \boxed{
 \frac1q\sum_{r\, (q)}|\widehat A_{C,q}(r)|^2
 =\sum_{n\in\Z}|A_C(n)|^2.}
 \label{eq:tpc32-no-wrap-parseval}
\end{equation}
If $C\le J$ and $C\ll_{\mathscr D}Q$, then, for every $\sigma>0$,
\begin{equation}
 \boxed{
 \#\left\{r\pmod q:
 |\widehat A_{C,q}(r)|>N_0X^{-\sigma}\right\}
 \ll_{\eps,h,\mathscr D}X^{2\sigma+\eps}.}
 \label{eq:tpc32-no-wrap-exceptional-count}
\end{equation}
The exceptional set in \eqref{eq:tpc32-no-wrap-exceptional-count} is
allowed to contain $r=0$.
\end{theorem}

\begin{proof}
Ordinary finite Fourier Parseval gives
\[
 \frac1q\sum_{r\, (q)}|\widehat A_{C,q}(r)|^2
 =\sum_{a\, (q)}
   \left|\sum_{n\equiv a\, (q)}A_C(n)\right|^2.
\]
The diameter condition in \eqref{eq:tpc32-no-wrap-prime} makes reduction
modulo $q$ injective on $\operatorname{supp}A_C$, so the last expression
is exactly $\sum_n|A_C(n)|^2$.  If $E$ frequencies exceed the threshold
in \eqref{eq:tpc32-no-wrap-exceptional-count}, then
\[
 E N_0^2X^{-2\sigma}
 \le\sum_{r\, (q)}|\widehat A_{C,q}(r)|^2
 =q\|A_C\|_2^2
 \ll X^\eps\frac qQ N_0^2.
\]
Now $q\asymp_{\mathscr D}Q$, and the result follows.
\end{proof}

The no-wrap choice is useful but not cosmetic.  If one merely takes an
arbitrary $q\asymp Q$, determinants separated by a multiple of $q$
occupy the same residue cell, and the right side of Parseval is the
energy of residue-cell sums rather than the coefficient energy itself.

\subsection{Farey-family dispersion}

For a reduced rational frequency $a/v$, put
\begin{equation}
 \widehat A_C(a/v)
 :=\sum_n A_C(n)e_v(-an).
 \label{eq:tpc32-rational-det-transform}
\end{equation}
The classical additive large sieve applies to an arbitrary complex
coefficient sequence; no factorability hypothesis is needed here
\citep{IwaniecKowalski2004}.

\begin{theorem}[Additive large sieve for determinant frequencies]
\label{thm:tpc32-farey-large-sieve}
For $V\ge2$, $C\le J$, and $C\ll_{\mathscr D}Q$,
\begin{equation}
 \boxed{
 \sum_{2\le v\le V}\ \sum_{a\, (v)}^{*}
 |\widehat A_C(a/v)|^2
 \ll_{\eps,h,\mathscr D}
 X^\eps(Q+V^2)\frac{N_0^2}{Q}.}
 \label{eq:tpc32-farey-large-sieve}
\end{equation}
Consequently
\begin{align}
 &\#\left\{(a,v):2\le v\le V, (a,v)=1,
 |\widehat A_C(a/v)|>N_0X^{-\sigma}\right\}
 \notag\\
 &\hspace{8em}\ll_{\eps,h,\mathscr D}
 X^{2\sigma+\eps}\left(1+\frac{V^2}{Q}\right).
 \label{eq:tpc32-farey-exceptional-count}
\end{align}
If $Q=X^{\beta+o(1)}$ and $V=X^{\rho+o(1)}$ with $\rho>0$, the
exceptional proportion among the $\asymp V^2$ reduced fractions is
power-small for every fixed
\begin{equation}
 \boxed{\sigma<\min\{\rho,\beta/2\}.}
 \label{eq:tpc32-farey-saving-range}
\end{equation}
\end{theorem}

\begin{proof}
Translate the integer support interval in
\eqref{eq:tpc32-determinant-support-interval} into an interval of
positive integers.  Translation changes each transform by a
unit-modulus factor.  The interval has length at most
$2B_{\mathscr D}Q+1$, so the additive large sieve gives
\[
 \sum_{2\le v\le V}\sum_{a\, (v)}^*
 |\widehat A_C(a/v)|^2
 \ll_{\mathscr D}(Q+V^2)\|A_C\|_2^2.
\]
Insert \eqref{eq:tpc32-det-l2-C-below-J}.  Chebyshev's inequality gives
\eqref{eq:tpc32-farey-exceptional-count}.

The number of reduced fractions with denominator at most $V$ is
$\asymp V^2$.  If $2\rho\le\beta$, division of
\eqref{eq:tpc32-farey-exceptional-count} by $V^2$ gives the exponent
$2\sigma-2\rho+o(1)$.  If $2\rho\ge\beta$, it gives
$2\sigma-\beta+o(1)$.  Both are negative precisely in the strict range
\eqref{eq:tpc32-farey-saving-range}.
\end{proof}

\begin{remark}[What the large sieve omits]
The sum in \eqref{eq:tpc32-farey-large-sieve} begins at $v=2$, so all
of its reduced frequencies are nonzero.  One may add the single point
$0/1$ to the large-sieve family, but the inequality then permits that
one value to be as large as the natural scale $N_0$.  A density-one
conclusion for rational frequencies therefore says nothing about the
distinguished coefficient $\widehat A_{C,q}(0)$ required by the
untwisted physical reassembly.  The next section isolates exactly the
extra input needed there.
\end{remark}
